\documentclass[11pt]{article}
\usepackage{amsmath, amssymb}
\usepackage{geometry}
\usepackage{parskip}

\geometry{margin=1.2in}

\title{Perikymata Distribution from the Final Curve}
\date{}

\begin{document}
\maketitle

Given the final outer enamel surface (OES) $\mathbf{r}_f(s) = (x_f(s), y_f(s))$ and the
termination time field $t^*(s)$, the perikymata distribution measures how much developmental
time is compressed into each spatial region of the OES. Perikymata are essential markesrs of enamel termination time, such that 1 perikymata corresponds to $\mathcal{R}$ days of development, where $\mathcal{R}$ is the Retzius periodicity. The following steps outline how to compute the perikymata distribution along the OES. In other words, 1 PK corresponds to the distance the termination front progression in $\mathcal{R}$ days.

%------------------------------------------------------------------
\section{Establish a Reference Axis}
%------------------------------------------------------------------

In general, the reference axis is specified as an arbitrary unit vector
$\hat{\mathbf{d}} \in \mathbb{R}^2$ and an origin point $\mathbf{r}_0$ on the OES.
The projected coordinate of each point along the crown is then
\begin{equation}
    \rho(s) = \bigl(\mathbf{r}_f(s) - \mathbf{r}_0\bigr) \cdot \hat{\mathbf{d}}.
\end{equation}
The integration bounds are $\rho \in [0, H]$ where $H = \max_s \rho(s)$. When the geometry
admits well-defined anatomical landmarks, $\hat{\mathbf{d}}$ and $\mathbf{r}_0$ can be set
naturally from the cusp and cervix; in the general case they are free parameters that define
the measurement orientation.

%------------------------------------------------------------------
\section{Partition the OES into Intervals}
%------------------------------------------------------------------

Divide $[0, H]$ into $N$ equal intervals of width $\Delta = H/N$. Define a smooth indicator
for the $k$-th interval:
\begin{equation}
    \chi_k(s)
    =
    \sigma\!\left(\beta\bigl(\rho(s) - k\Delta\bigr)\right)
    \cdot
    \sigma\!\left(\beta\bigl((k+1)\Delta - \rho(s)\bigr)\right),
\end{equation}
where $\sigma$ is the logistic function and $\beta$ controls sharpness.

%------------------------------------------------------------------
\section{Integrate Developmental Time per Interval}
%------------------------------------------------------------------

\begin{equation}
    \mathcal{T}_k
    =
    \int_0^L \left|\frac{dt^*}{ds}\right| \chi_k(s)\, ds
    \;\approx\;
    \sum_i dt^*_i \cdot \chi_k(s_{i+1/2}),
\end{equation}
where $dt^*_i = |t^*(s_{i+1}) - t^*(s_i)|$.

%------------------------------------------------------------------
\section{Convert to Perikymata Count}
%------------------------------------------------------------------

Given Retzius periodicity $\mathcal{R}$ (days):
\begin{equation}
    P_k = \frac{365\,\mathcal{T}_k}{\mathcal{R}}.
\end{equation}
Non-uniform $t^*(s)$ directly produces non-uniform $P_k$, since regions where the off-front
decelerates accumulate more $dt^*$ per unit projected length.

%------------------------------------------------------------------
\section{True PK distribution}
%------------------------------------------------------------------

In this regime, calculate number of perikymata as a function of distance along the arc length $S$ of the OES (as opposed to the $s$ arclength we get from the EDJ). Ultimatly, one of hte things we want to understand is how $S = S(s)$, that is the mapping from EDJ arclength to OES arclength. 
\end{document}
